\documentclass[../../main.tex]{subfiles}
\begin{document}

\begin{flushright}
  \footnotesize\textit{【自動文字起こし・要確認】}
\end{flushright}

\noindent
[\textbf{解}] 3点 $A(\cos\alpha, \sin\alpha), B(\cos\beta, \sin\beta), C(\cos\gamma, \sin\gamma)$ を頂点とする $\Delta ABC$ を考える. この外心は $(0,0)$ で, 題意から重心も $(0,0)$ である.

外心と重心が一致するのは正三角形なので, $\Delta ABC$ は正三角形.

(以上により $0 \le \alpha < \beta < \gamma \le 2\pi$ を用いた).

従って
\begin{align}
  \beta - \alpha = \gamma - \beta = \dfrac{2\pi}{3} \quad /\!/
\end{align}

\bigskip
\noindent
[\textbf{解2}]
$\cos^2\gamma + \sin^2\gamma = 1$ に代入
\begin{align}
  2 + 2\cos(\beta - \alpha) = 1
\end{align}
\begin{align}
  \cos(\beta - \alpha) = -\dfrac{1}{2}
\end{align}
\begin{align}
  \therefore \beta - \alpha = \dfrac{2}{3}\pi, \, \dfrac{4}{3}\pi
\end{align}
同様にして
\begin{align}
  \gamma - \alpha = \dfrac{2}{3}\pi, \, \dfrac{4}{3}\pi
\end{align}
以上と $0 \le \alpha < \beta < \gamma < 2\pi$ から
\begin{align}
  \beta - \alpha = \dfrac{2}{3}\pi, \quad \gamma - \alpha = \dfrac{4}{3}\pi
\end{align}
\begin{align}
  \therefore \gamma - \beta = \dfrac{2}{3}\pi \quad /\!/
\end{align}

\newpage

\end{document}
